There is a finite extension @@M0@@ of @@M1@@ such that @@M2@@ has all its residue field extensions trivial. For an integral scheme @@M3@@ finite and flat over an open subset @@M4@@ of @@M5@@, with function field @@M6@@, there is an open subset @@M7@@ of @@M8@@ such that
@@M9@@
over @@M10@@; after shrinking @@M11@@ and @@M12@@, we may suppose that @@M13@@ and @@M14@@. Since the assertion to be proved is local for f.p.q.c., we may replace @@M15@@ by @@M16@@. Finally, apply 3.1.

\sgaStableTarget{sga6.e12.rem.3-6}\textbf{3.6. Remark.} It follows in the same way, taking into account the exact sequences @@M17@@, that, for every @@M18@@,
@@M19@@
is a unipotent group.

\sgaStableTarget{sga6.e12.sec.4}\section*{4. Proof of II*: the most delicate part of the proof}

Recall assertion II*: Theorem 1.1 when @@M20@@ and @@M21@@ are integral over @@M22@@ and the sequence (*) @@M23@@ is exact; applying the generic flatness theorem EGA IV 6.9.1 and EGA III 7.8.8 reduces us to proving the

\sgaStableTarget{sga6.e12.prop.4-1}\textbf{Proposition 4.1.} \uline{Let} @@M24@@ \uline{be an integral noetherian scheme,} @@M25@@ \uline{and} @@M26@@ \uline{proper morphisms of finite presentation,} @@M27@@ \uline{a surjective} @@M28@@\uline{-morphism, and} @@M29@@ \uline{the canonical morphism. Suppose that}
\begin{enumerate}[label=\arabic*)]
\item @@M30@@ \uline{are flat over} @@M31@@.
\item @@M32@@, @@M33@@ \uline{universally.}

